\chapter{COMPLEXITY ANALYSIS AND EMPIRICAL VALIDATION}\label{chap:complexity}

This chapter presents the formal time and space complexity of both algorithms and validates the theoretical bounds with actual experimental measurements.

\section{Analytical Methodology}

Complexity was analyzed using the standard \textbf{operation-counting} approach~\citep{cormen2009}: (1) identify all operations (heap insertions, extractions, edge relaxations), (2) count the maximum number of each as a function of $|V|$ and $|E|$, (3) multiply by per-operation cost, and (4) verify empirically.

\section{Time Complexity}

\subsection{Baseline Dijkstra}

\begin{enumerate}
    \item \textbf{Initialization:} Setting $\mathit{dist}[v]=\infty$ and $\mathit{parent}[v]=\texttt{NIL}$ for all $v$ costs $O(|V|)$.
    \item \textbf{Priority queue insertions:} Each node is inserted initially once and re-inserted (lazy) at most once per incoming edge. Total insertions $\leq |V|+|E|$. Each costs $O(\log|V|)$ (binary heap sift-up). Total: $O(|E|\log|V|)$.
    \item \textbf{Extract-min operations:} Total extractions $\leq |V|+|E|$, each $O(\log|V|)$. Total: $O(|E|\log|V|)$.
    \item \textbf{Edge relaxations:} For each settled node $u$, all $\deg(u)$ outgoing edges are examined. Across the full algorithm: $\sum_{u\in V}\deg(u)=|E|$. Each relaxation costs $O(1)$ (comparison + hash map lookup). Total: $O(|E|)$.
\end{enumerate}

\[
T_{\text{baseline}} = O(|E|\log|V|)
\]

In sparse graphs where $|E|=O(|V|)$, this simplifies to $O(|V|\log|V|)$.

\textit{Note:} A Fibonacci heap would yield $O(|E|+|V|\log|V|)$~\citep{fredman1987fibonacci}, but binary heaps provide better constant factors in practice.

\subsection{Probability-Pruned Dijkstra}

\begin{enumerate}
    \item \textbf{Worst case ($\tau=0$):} Identical to the baseline: $O(|E|\log|V|)$.
    \item \textbf{Empirically observed case ($\tau>0$):} Let $|E'|\leq|E|$ denote edges surviving the threshold. Pruned edges are checked in $O(1)$ without heap insertion. The observed complexity is:
    \[
    T_{\text{pruned}} = O(|E'|\log|V|),\quad |E'|\leq|E|
    \]
    This bound is not formally proven in the worst case---$|E'|$ depends on the input graph's probability distribution and the chosen $\tau$---but is consistently confirmed by the empirical measurements in Section~\ref{sec:empirical_time}.
    \item \textbf{Pruning ratio:} Define $\rho = 1 - |E'|/|E|$. Higher $\tau$ increases $\rho$, reducing runtime.
\end{enumerate}

\subsection{Empirical Time Validation}
\label{sec:empirical_time}

Table~\ref{table:time_by_size} shows measured wall-clock times at $\tau=0.01$ across graph sizes.

\begin{table}[htbp]
\centering
\renewcommand{\arraystretch}{1.3}
\begin{tabular}{@{}rrrrr@{}}
\toprule
$|V|$ & \textbf{Baseline (ms)} & \textbf{Pruned (ms)} & \textbf{Wall Speedup} & \textbf{Edge Speedup} \\
\midrule
1,000  & 1.67 & 0.189 & 9.4$\times$ & 31.5$\times$ \\
5,000  & 7.12 & 0.215 & 43.1$\times$ & 104.0$\times$ \\
10,000 & 15.83 & 0.222 & 77.0$\times$ & 180.4$\times$ \\
\bottomrule
\end{tabular}
\caption[Execution time scaling by graph size]{Median execution time at $\tau=0.01$ by graph size.}
\label{table:time_by_size}
\end{table}

\noindent The baseline times scale as $1.67 \to 7.12 \to 15.83$ ms for $|V|=1{,}000 \to 5{,}000 \to 10{,}000$, closely matching the $O(|V|\log|V|)$ prediction for sparse graphs ($|E|=O(|V|)$). The pruned algorithm's nearly constant time ($\sim$0.2\,ms) is consistent with the hypothesis that $|E'|$ grows much slower than $|E|$ as the graph scales, though only three data points are available, so a definitive growth rate cannot be claimed.

\section{Space Complexity}

\begin{table}[htbp]
\centering
\renewcommand{\arraystretch}{1.3}
\begin{tabular}{@{}lll@{}}
\toprule
\textbf{Data Structure} & \textbf{Space} & \textbf{Purpose} \\
\midrule
Adjacency list & $O(|V|+|E|)$ & Graph representation \\
Distance map $\mathit{dist}[v]$ & $O(|V|)$ & Shortest path estimates \\
Parent map $\mathit{parent}[v]$ & $O(|V|)$ & Path reconstruction \\
Priority queue (heap) & $O(|V|+|E|)$ worst case & Node ordering \\
\midrule
\textbf{Total} & $\mathbf{O(|V|+|E|)}$ & \\
\bottomrule
\end{tabular}
\caption[Space complexity breakdown]{Space complexity breakdown.}
\label{table:space_complexity}
\end{table}

The pruning threshold $\tau$ adds no asymptotic space overhead. In practice, the pruned algorithm uses less peak memory because fewer heap entries are created.

\subsection{Empirical Memory Validation}

\begin{table}[htbp]
\centering
\renewcommand{\arraystretch}{1.3}
\begin{tabular}{@{}rrrr@{}}
\toprule
$|V|$ & \textbf{Baseline Peak (KB)} & \textbf{Pruned Peak (KB)} & \textbf{Reduction} \\
\midrule
1,000  & 83.8 & 4.4 & 19.0$\times$ \\
5,000  & 342.9 & 4.5 & 75.9$\times$ \\
10,000 & 655.0 & 4.5 & 145.9$\times$ \\
\bottomrule
\end{tabular}
\caption[Peak memory scaling by graph size]{Peak memory at $\tau=0.01$ by graph size.}
\label{table:memory_by_size}
\end{table}

\noindent Peak \texttt{tracemalloc} allocation measures algorithm data structures only (distance dict, parent dict, priority-queue heap); the graph adjacency list is built before tracing begins and is excluded. Baseline peak memory (83.8~KB $\to$ 342.9~KB $\to$ 655.0~KB) scales linearly with $|V|$, confirming the $O(|V|+|E|)$ bound, while the pruned algorithm's constant $\approx$4.5~KB at all sizes demonstrates that only a small frontier is ever materialized.

\section{Summary of Complexity}

\begin{table}[htbp]
\centering
\renewcommand{\arraystretch}{1.4}
\small
\begin{tabular}{@{}lccc@{}}
\toprule
\textbf{Algorithm} & \textbf{Time (worst)} & \textbf{Time (empirical)} & \textbf{Space} \\
\midrule
Baseline Dijkstra & $O(|E|\log|V|)$ & $O(|E|\log|V|)$ & $O(|V|+|E|)$ \\
Pruned ($\tau=0$) & $O(|E|\log|V|)$ & $O(|E|\log|V|)$ & $O(|V|+|E|)$ \\
Pruned ($\tau>0$) & $O(|E|\log|V|)$ & $O(|E'|\log|V|)$ & $O(|V|+|E|)$ \\
\bottomrule
\end{tabular}
\caption[Complexity summary]{Complexity summary.}
\label{table:complexity_summary}
\end{table}

\noindent Pruning reduces empirically observed time without affecting worst-case or space bounds. The empirical bound $O(|E'|\log|V|)$ is consistently confirmed by experiment (Section~\ref{sec:empirical_time}) but depends on the input distribution.

\noindent\textbf{Key empirical findings:}
Baseline wall-clock time scales as $O(|V|\log|V|)$ for sparse graphs, matching theory. Pruned wall-clock time is nearly constant across graph sizes at fixed $\tau$, consistent with $|E'| \ll |E|$. Peak memory scales linearly for the baseline, while pruned memory is effectively constant. Edge-relaxation speedups (up to 1,986$\times$) exceed wall-clock speedups (up to 738$\times$), the difference attributable to fixed Python overhead per function call.
